\section{Equity Analysis}

\subsection{C3 Alignment: Freight Corridors and Economic Opportunity}

The ROUTE rubric's C3 dimension (Economic Opportunity Access) scores interstate corridors on the GDP per capita of the counties within 20 miles of the corridor alignment, relative to the national median. A C3 score of 5.0 corresponds to a corridor buffer with GDP per capita equal to the national median; scores above 5.0 indicate below-median GDP (economically disadvantaged), and scores below 5.0 indicate above-median GDP (economically advantaged) \citep{ROUTE_A1}.

The T1/T1 hub locations are determined by freight intensity and network connectivity, not by equity criteria. Yet when we compute the C3 score of the 30-mile buffer around each T1/T1 hub location, 7 of the 9 hubs have buffer C3 scores exceeding 5.0 --- meaning their surrounding communities have below-average GDP per capita. The correlation between C3 score at the hub location and transit-dependent household density in the hub 30-mile buffer is $r = 0.68$ (Pearson, $n = 9$ hubs, $p < 0.05$). The rubric-driven site selection for freight infrastructure produces, as a byproduct, near-optimal siting for transit equity.

\begin{table}[h]
\centering
\caption{C3 Score and Transit-Dependent Household Density by T1/T1 Hub}
\label{tab:equity}
\begin{tabular}{lrr}
\toprule
Hub Location & C3 Score (Buffer) & Transit-Dep. HH Density (per sq mi) \\
\midrule
Gary/Chicago, IL & 7.2 & 142 \\
San Antonio, TX & 6.8 & 118 \\
Toledo, OH & 6.4 & 89 \\
Atlanta, GA & 6.1 & 96 \\
Jacksonville, FL & 5.9 & 78 \\
Richmond, VA & 5.6 & 71 \\
Sacramento, CA & 5.1 & 58 \\
Seattle, WA & 4.3 & 62 \\
Boston, MA & 3.8 & 48 \\
\midrule
Pearson $r$ (C3 vs density) & \multicolumn{2}{c}{0.68} \\
\bottomrule
\end{tabular}
\end{table}

\subsection{Community-Level Equity Analysis}

Three hub locations warrant specific attention for their equity implications:

\textbf{Gary/Chicago south side (I-80/I-90)}: The Gary/Chicago hub has the highest transit-dependent household density of any T1/T1 location. Gary, Indiana --- a post-industrial city with above-average poverty and below-average vehicle ownership --- is the primary beneficiary of the hub, not Chicago's existing transit system. Chicago's south side communities (south of I-290) are within 20 miles of the Gary hub, accessible by bus or bicycle, and have significantly below-average intercity transit access despite their proximity to a major metropolitan area. ACS data shows zero-vehicle household rates of 28--34\% in Gary census tracts, compared to a national average of 8.7\% \citep{ACS_B08201_2022}. The Gary hub serves a transit-dependent population that has been systematically underserved by both rail (no Amtrak station in Gary proper) and bus (Greyhound Gary terminal closed 2021).

\textbf{San Antonio south side (I-10/I-35)}: San Antonio's T1/T1 hub is located in the southwest quadrant of the city near the I-10/I-35 interchange, which is geographically closer to the primarily Hispanic communities of south and west San Antonio than to the downtown Amtrak/Greyhound terminal. ACS data for San Antonio's southern census tracts (Bexar County census tracts with median household income below \$35,000) shows zero-vehicle household rates of 18--24\% \citep{ACS_B08201_2022}, compared to 8.7\% nationally. These communities currently rely on local VIA Metropolitan Transit for intra-city trips; there is no viable intercity transit connection for Laredo, Corpus Christi, or Houston without a personal vehicle.

\textbf{Rural T1/T2 stops along I-69 and I-29S}: The 12 T1/T2 stops along the I-69 alignment (Texarkana TX to Indianapolis IN) and I-29S alignment (Laredo TX to Sioux Falls SD) serve agricultural worker communities that are among the most transit-isolated in the country. These communities are not near any T1/T1 hub; they depend on the T1/T2 stop network for any intercity transit connection. ACS data shows zero-vehicle household rates of 12--18\% in the counties along these alignments \citep{ACS_B08201_2022}, with high concentrations of Hispanic and Native American households and below-average English language proficiency --- populations for whom bus travel is often the preferred or only viable intercity mode.

\subsection{The Equity Multiplier Effect}

The core equity argument is that the I2.0 highway investment has a transit equity dividend that is not captured in the freight benefit-cost analysis. The transit layer adds \$2 billion to a \$253 billion program and serves 12.4 million transit-dependent Americans. But the equity value of connecting these communities is not captured in the shipper reliability or freight time savings metrics that dominate the freight B/C calculation.

The equity multiplier operates through three channels:
\begin{enumerate}
  \item \textbf{Access to labor markets}: intercity bus connectivity allows workers in small cities and rural communities to access employment in larger metros for interviews, training, and temporary work without car ownership
  \item \textbf{Access to healthcare}: studies of rural health access consistently identify lack of transportation as the primary barrier to specialty care for rural populations; intercity bus connectivity to hub cities with medical facilities addresses this gap \citep{Pisarski2006}
  \item \textbf{Access to family}: transit-dependent populations with dispersed family networks --- a structural characteristic of low-income communities with high mobility --- rely on intercity bus for family visits and emergency travel
\end{enumerate}

\subsection{Recommended Equity Standard}

This paper recommends that T1 hub transit service be designated as Essential Intercity Transportation under the Surface Transportation Block Grant program, with the following minimum service requirements:
\begin{itemize}
  \item Minimum 2 round trips per day between each T1/T1 hub and the T1/T2 regional stops within 120 miles
  \item Regulated maximum fares, set at not more than 1.5 times the per-mile Amtrak unreserved coach fare for comparable distances
  \item ADA-accessible vehicles on all T1 hub routes
  \item Advance reservation system with online and telephone booking
\end{itemize}

Without service frequency and fare requirements, the hub infrastructure will not produce the equity outcomes its location and cost efficiency suggest are achievable. Private bus operators will serve routes with proven demand and acceptable load factors; they will not voluntarily serve routes that serve primarily transit-dependent, low-income populations at the frequency those populations require. The Essential Intercity Transportation designation provides the regulatory basis for requiring minimum service standards as a condition of hub access agreements.
